\documentclass[10pt,a4paper]{amsart}
\usepackage{amsmath,amssymb,amsthm,graphicx}
% Retained minimal preamble for the Stacks Project chapter more-algebra.tex (More on Algebra).
% The chapter file pinned-original-source.tex carries only \input{preamble} on its first line, so no
% preamble text lives inside the chapter itself.  The Stacks Project uses one single shared preamble for
% every chapter; the snapshot stages it identically as sources/stacks-chow/preamble.tex and
% sources/stacks-etale-cohomology/preamble.tex.  The block below is copied mechanically, line for line,
% from sources/stacks-chow/preamble.tex lines 14-18 (the Xy-pic setup) and 154-217 (the native theorem
% declarations together with the mathematical macros).  Nothing was added, removed or redefined; the
% wrapper class and the amsmath/amssymb/amsthm/graphicx base packages are supplied by the assembler.
% Chapter: more-algebra.tex -- More on Algebra.
% For commutative diagrams we use Xy-pic
\usepackage[all]{xy}

% We use 2cell for 2-commutative diagrams.
\xyoption{2cell}

% Theorem environments.
%
\theoremstyle{plain}
\newtheorem{theorem}[subsection]{Theorem}
\newtheorem{proposition}[subsection]{Proposition}
\newtheorem{lemma}[subsection]{Lemma}

\theoremstyle{definition}
\newtheorem{definition}[subsection]{Definition}
\newtheorem{example}[subsection]{Example}
\newtheorem{exercise}[subsection]{Exercise}
\newtheorem{situation}[subsection]{Situation}

\theoremstyle{remark}
\newtheorem{remark}[subsection]{Remark}
\newtheorem{remarks}[subsection]{Remarks}

\numberwithin{equation}{subsection}

% Macros
%
\def\lim{\mathop{\mathrm{lim}}\nolimits}
\def\colim{\mathop{\mathrm{colim}}\nolimits}
\def\Spec{\mathop{\mathrm{Spec}}}
\def\Hom{\mathop{\mathrm{Hom}}\nolimits}
\def\Ext{\mathop{\mathrm{Ext}}\nolimits}
\def\SheafHom{\mathop{\mathcal{H}\!\mathit{om}}\nolimits}
\def\SheafExt{\mathop{\mathcal{E}\!\mathit{xt}}\nolimits}
\def\Sch{\mathit{Sch}}
\def\Mor{\mathop{\mathrm{Mor}}\nolimits}
\def\Ob{\mathop{\mathrm{Ob}}\nolimits}
\def\Sh{\mathop{\mathit{Sh}}\nolimits}
\def\NL{\mathop{N\!L}\nolimits}
\def\CH{\mathop{\mathrm{CH}}\nolimits}
\def\proetale{{pro\text{-}\acute{e}tale}}
\def\etale{{\acute{e}tale}}
\def\QCoh{\mathit{QCoh}}
\def\Ker{\mathop{\mathrm{Ker}}}
\def\Im{\mathop{\mathrm{Im}}}
\def\Coker{\mathop{\mathrm{Coker}}}
\def\Coim{\mathop{\mathrm{Coim}}}

% Boxtimes
%
\DeclareMathSymbol{\boxtimes}{\mathbin}{AMSa}{"02}

%
% Macros for moduli stacks/spaces
%
\def\QCohstack{\mathcal{QC}\!\mathit{oh}}
\def\Cohstack{\mathcal{C}\!\mathit{oh}}
\def\Spacesstack{\mathcal{S}\!\mathit{paces}}
\def\Quotfunctor{\mathrm{Quot}}
\def\Hilbfunctor{\mathrm{Hilb}}
\def\Curvesstack{\mathcal{C}\!\mathit{urves}}
\def\Polarizedstack{\mathcal{P}\!\mathit{olarized}}
\def\Complexesstack{\mathcal{C}\!\mathit{omplexes}}
% \Pic is the operator that assigns to X its picard group, usage \Pic(X)
% \Picardstack_{X/B} denotes the Picard stack of X over B
% \Picardfunctor_{X/B} denotes the Picard functor of X over B
\def\Pic{\mathop{\mathrm{Pic}}\nolimits}
\def\Picardstack{\mathcal{P}\!\mathit{ic}}
\def\Picardfunctor{\mathrm{Pic}}
\def\Deformationcategory{\mathcal{D}\!\mathit{ef}}

% --- appended 2026-10-12 (Chief): the project own definitions of the reference/slogan/history
% display environments, copied mechanically from the shared Stacks preamble lines 8-12
% (\usepackage{verbatim} plus the three \newenvironment lines, which the chapter preamble keeps
% verbatim from the project preamble). Nothing added, removed or redefined: these environments
% carry no mathematical content and the project itself typesets them as comments. Without these
% three lines a counted statement whose own text uses them fails with Environment slogan undefined.
\usepackage{verbatim}
\newenvironment{reference}{\comment}{\endcomment}
\newenvironment{slogan}{\comment}{\endcomment}
\newenvironment{history}{\comment}{\endcomment}
\usepackage{cleveref}
\crefname{originallemma}{original lemma}{original lemmas}
\Crefname{originallemma}{Original lemma}{Original lemmas}
\makeatletter
\newlabel{lemma-largest-ideal-henselian}{{0G1S}{source}{}{sourcepointer-1}{}}
\newlabel{lemma-largest-ideal-henselian@cref}{{[originallemma][0][]0G1S}{[1][0][]source}{}{}{}}
\makeatother
\crefname{originaldefinition}{original definition}{original definitions}
\Crefname{originaldefinition}{Original definition}{Original definitions}
\makeatletter
\newlabel{definition-derived-complete}{{091S}{source}{}{sourcepointer-2}{}}
\newlabel{definition-derived-complete@cref}{{[originaldefinition][0][]091S}{[1][0][]source}{}{}{}}
\makeatother
\crefname{originallemma}{original lemma}{original lemmas}
\Crefname{originallemma}{Original lemma}{Original lemmas}
\makeatletter
\newlabel{lemma-complete-derived-complete}{{091R}{source}{}{sourcepointer-3}{}}
\newlabel{lemma-complete-derived-complete@cref}{{[originallemma][0][]091R}{[1][0][]source}{}{}{}}
\makeatother
\crefname{originallemma}{original lemma}{original lemmas}
\Crefname{originallemma}{Original lemma}{Original lemmas}
\makeatletter
\newlabel{lemma-complete-henselian}{{0ALJ}{source}{}{sourcepointer-4}{}}
\newlabel{lemma-complete-henselian@cref}{{[originallemma][0][]0ALJ}{[1][0][]source}{}{}{}}
\makeatother
\crefname{originallemma}{original lemma}{original lemmas}
\Crefname{originallemma}{Original lemma}{Original lemmas}
\makeatletter
\newlabel{lemma-henselian-henselian-pair}{{0DYD}{source}{}{sourcepointer-5}{}}
\newlabel{lemma-henselian-henselian-pair@cref}{{[originallemma][0][]0DYD}{[1][0][]source}{}{}{}}
\makeatother
\crefname{originallemma}{original lemma}{original lemmas}
\Crefname{originallemma}{Original lemma}{Original lemmas}
\makeatletter
\newlabel{lemma-kernel-to-completion-square-zero}{{0G3G}{source}{}{sourcepointer-6}{}}
\newlabel{lemma-kernel-to-completion-square-zero@cref}{{[originallemma][0][]0G3G}{[1][0][]source}{}{}{}}
\makeatother
\crefname{originallemma}{original lemma}{original lemmas}
\Crefname{originallemma}{Original lemma}{Original lemmas}
\makeatletter
\newlabel{lemma-locally-nilpotent-henselian}{{0ALI}{source}{}{sourcepointer-7}{}}
\newlabel{lemma-locally-nilpotent-henselian@cref}{{[originallemma][0][]0ALI}{[1][0][]source}{}{}{}}
\makeatother
\title{A ring derived complete with respect to an ideal forms a henselian pair}
\author{The Stacks Project Authors}
\date{stacks-project:more-algebra.tex:local-snapshot-20261012; GNU FDL 1.2 or later}
\begin{document}
\maketitle
\setcounter{subsection}{1}
\noindent\textit{Editorial note.} All delivered mathematical text is the author's own, copied from the pinned Stacks Project chapter file more-algebra.tex; no statement or proof step was rewritten, reordered or removed. The counted span is the author environment and its own proof environment. References inside the span are handed to the assembler as navigation stubs that carry the canonical tag of the cited tag only; whether each of them is load-bearing is recorded on the stub itself.\par\medskip
\noindent\textit{Allowed standard prerequisites (editorial).}\begin{itemize}
\item A pair is henselian as soon as it is henselian with respect to every principal ideal generated by an element of the ideal. invoked at proof lines 27092-27094. Reduces the statement for the given ideal to the case of a single element.
\item A complete ring with respect to a principal ideal forms a henselian pair, and a pair is henselian when its completion is and the kernel of the completion map is handled. invoked at proof lines 27096-27102. Transfers the property from the adic completion back to the ring through the two intermediate statements.
\item A ring that is derived complete with respect to an element is derived complete, hence complete, with respect to the principal ideal that element generates. used at proof lines 27094-27096. Justifies applying the completeness statements to the principal ideal generated by the chosen element.
\end{itemize}\medskip
\ifdefined\hypertarget\hypertarget{sourcepointer-1}{}\fi
\noindent\textit{Source reference.} Original lemma 0G1S: sources/stacks-more-algebra/pinned-original-source.tex, source lines 2938--2938. Invoked at proof lines 27092-27094 to reduce the statement to the ideal generated by a single element; every step after the reduction is inside the counted span.\par\medskip
\ifdefined\hypertarget\hypertarget{sourcepointer-2}{}\fi
\noindent\textit{Source reference.} Original definition 091S: sources/stacks-more-algebra/pinned-original-source.tex, source lines 26021--26021. Cited at proof lines 27094-27096 for the immediate consequence that the ring is complete with respect to the principal ideal; the use made of that completeness is inside the counted span.\par\medskip
\ifdefined\hypertarget\hypertarget{sourcepointer-3}{}\fi
\noindent\textit{Source reference.} Original lemma 091R: sources/stacks-more-algebra/pinned-original-source.tex, source lines 25963--25963. Used at proof lines 27096-27098 to see that the map from the ring to its completion is surjective; the transfer of henselianity that follows is inside the counted span.\par\medskip
\ifdefined\hypertarget\hypertarget{sourcepointer-4}{}\fi
\noindent\textit{Source reference.} Original lemma 0ALJ: sources/stacks-more-algebra/pinned-original-source.tex, source lines 2483--2483. Invoked at proof lines 27098-27100 to see that the pair with the completion is henselian; the reduction to the intersection of the powers that uses it is inside the counted span.\par\medskip
\ifdefined\hypertarget\hypertarget{sourcepointer-5}{}\fi
\noindent\textit{Source reference.} Original lemma 0DYD: sources/stacks-more-algebra/pinned-original-source.tex, source lines 2766--2766. Cited at proof lines 27100-27102 to reduce to the pair formed with the intersection of the powers of the ideal; the final two references that settle that case are inside the counted span.\par\medskip
\ifdefined\hypertarget\hypertarget{sourcepointer-6}{}\fi
\noindent\textit{Source reference.} Original lemma 0G3G: sources/stacks-more-algebra/pinned-original-source.tex, source lines 27068--27068. One of the two statements invoked at proof lines 27102-27103 to finish; the reduction stage that makes them applicable is inside the counted span.\par\medskip
\ifdefined\hypertarget\hypertarget{sourcepointer-7}{}\fi
\noindent\textit{Source reference.} Original lemma 0ALI: sources/stacks-more-algebra/pinned-original-source.tex, source lines 2417--2417. The second of the two statements invoked at proof lines 27102-27103 to finish the proof; the preceding reductions are inside the counted span.\par\medskip
\begin{lemma}
\label{lemma-derived-complete-henselian}
Let $A$ be a ring derived complete with respect to an ideal $I$.
Then $(A, I)$ is a henselian pair.
\end{lemma}


\begin{definition}
\label{definition-henselian-pair}
A {\it henselian pair} is a pair $(A, I)$ satisfying
\begin{enumerate}
\item $I$ is contained in the Jacobson radical of $A$, and
\item for any monic polynomial $f \in A[T]$ and factorization
$\overline{f} = g_0h_0$ with $g_0, h_0 \in A/I[T]$ monic
generating the unit ideal in $A/I[T]$, there
exists a factorization $f = gh$ in $A[T]$ with $g, h$ monic
and $g_0 = \overline{g}$ and $h_0 = \overline{h}$.
\end{enumerate}
\end{definition}

\begin{proof}
Let $f \in I$. By Lemma \ref{lemma-largest-ideal-henselian}
it suffices to show that $(A, fA)$ is a henselian pair.
Observe that $A$ is derived complete with respect to $fA$
(follows immediately from Definition \ref{definition-derived-complete}).
By Lemma \ref{lemma-complete-derived-complete} the map from
$A$ to the $f$-adic completion $A'$ of $A$ is surjective.
By Lemma \ref{lemma-complete-henselian} the pair $(A', fA')$
is henselian. Thus it suffices to show that $(A, \bigcap f^nA)$
is a henselian pair, see Lemma \ref{lemma-henselian-henselian-pair}.
This follows from Lemmas \ref{lemma-kernel-to-completion-square-zero}
and \ref{lemma-locally-nilpotent-henselian}.
\end{proof}

\end{document}
